
\section{Related Work}

\subsection{Linear Attention}

\textbf{Linear Transformer} (Katharopoulos et al., 2020). Uses $\phi(x) = \text{elu}(x) + 1$ 
without positional reweighting. First to show $\mathcal{O}(L)$ attention 
with reasonable performance.

\textbf{Performer} (Choromanski et al., 2021). Uses random Fourier features 
(RFF) for kernel approximation with theoretical guarantees.

\textbf{Nystr mformer} (Xiong et al., 2021). Uses Nystr m landmarks for low-rank 
attention approximation.

\subsection{Positional Reweighting}

\textbf{Cosformer} (Qin et al., 2022). The direct inspiration for this work. 
Uses cosine/sine reweighting to preserve locality.

\textbf{RoPE} (Su et al., 2021). Rotary position embeddings, standard in modern 
decoders. Not the same as Cosformer's reweighting -- RoPE rotates features, 
Cosformer scales them.

\textbf{ALiBi} (Press et al., 2022). Linear bias on attention scores. Preserves 
locality.

\subsection{Causal Linear Attention}

For causal models, naive linear attention is not enough -- the denominator 
must accumulate causally. Two main approaches:

\textbf{Cumulative sum (cumsum).} Compute running sums over the sequence. 
$\mathcal{O}(L)$ memory, single kernel.

\textbf{Chunked prefix-sum.} Split sequence into chunks; compute intra-chunk 
contributions directly, inter-chunk via state. Used in Cosformer, DeltaNet, 
and others.

\subsection{Pretrained Weight Transfer}

\textbf{DistilBERT} (Sanh et al., 2019). Distillation with weight transfer.

\textbf{LoRA} (Hu et al., 2021). Preserves pretrained weights, adds low-rank 
adapters.

\textbf{Our approach.} Combines exact weight transfer (Q/K/V/O) with LoRA recovery.

\subsection{Position of Our Work}

\begin{table}[H]
\centering
\small
\begin{tabular}{lll}
\toprule
\textbf{Axis} & \textbf{Prior work} & \textbf{Ours} \\
\midrule
\textbf{Method} & Cosformer (published) & Port to pretrained \\
\textbf{Weight transfer} & Rare & Exact, four architectures \\
\textbf{Bidirectional} & Sometimes & Dedicated path \\
\textbf{RoPE support} & Limited & Full support \\
\textbf{Recovery} & From scratch & LoRA \\
\bottomrule
\end{tabular}
\end{table}

\textbf{Key distinction.} Unlike most linear-attention papers that train from 
scratch, we focus on \textbf{weight-preserving port} to pretrained models. This 
requires careful handling of the \texttt{copy\_weights} procedure across 
architectures.



\subsection{Comparison Table}

\begin{lstlisting}

---
\end{lstlisting}